This project set out to investigate how embedded systems and fall detection algorithms can be used to support elderly citizens living independently by detecting falls and ensuring timely notification of relatives. Through analysis of demographic trends, existing solutions, and fall detection algorithms, a clear need was identified for a reliable, user-friendly, and economically accessible fall detection system.

Based on insights from the state of art analysis and an interview with a representative from the target group, a set of functional and non-functional requirements was defined. These requirements guided the design of the Lapsus system, which consists of a wrist-worn embedded device, a mobile application, and a cloud-based backend. The system is designed to automatically detect falls using accelerometer-based sensing and to notify relatives directly via \gls{sms}, ensuring fast and reliable communication without reliance on subscription-based services or internal alarm centers.

The implemented prototype demonstrates that threshold-based fall detection algorithms can be effectively deployed on low-power wearable hardware while maintaining acceptable accuracy and responsiveness. The wrist-mounted form factor, although not technically optimal compared to waist-mounted solutions, was chosen as a conscious trade-off to improve comfort, familiarity, and user acceptance among elderly users. Testing results show that the system is capable of detecting falls and delivering notifications within the defined constraints, thereby fulfilling the core requirements derived from the problem statement.

Furthermore, the project highlights important trade-offs between usability, cost, privacy, and security. Design choices such as phone number-based authentication, continuous \gls{gps} collection, and \gls{sms}-based alerts prioritize simplicity and reliability but also introduce potential privacy and security concerns. These limitations were acknowledged and discussed, emphasizing that additional safeguards and data management mechanisms would be necessary for real-world deployment.

In conclusion, the project demonstrates that embedded systems combined with fall detection algorithms can meaningfully support elderly citizens living independently by enhancing safety and enabling timely assistance from relatives. While the developed prototype is not production-ready, it successfully validates the technical feasibility of the proposed solution and identifies clear directions for future work, including improved data handling, enhanced security measures, and extended validation with real users.